Itt kezdődik a bevezetés.\par
Lehetőség szerint a különböző szekciókat külön fájlokba írjuk, az átláthatóság érdekében, majd az \texttt{\textbackslash input\{./sections/beveztés.tex\}} paranccsal importáljuk őket. Ezek a fájlok úgy viselkednek, mint a sztendern megszokott \texttt{.tex} fájlok, viszont nem kell az elejére preambulum, mivel örökli ezt a \texttt{main.tex} fájlból.

\subsection{Dolgozat tagolása}

A dolgozatot érdemes részekre tagolni, amiknek rövid leíró címeket lehet adni.

\subsubsection{Részletesebb tagolás}

Amennyiben szükséges, az alpontokat is tovább lehet bontani a \texttt{\textbackslash subsubsection\{\}} paranccsal.\par

\subsection{Tartalomjegyzék megjelenés}

Ezek a tagolások megjelennek a tartalomjegyzékben is, mint klikkelhető linkek!
